\input{preamble}

\begin{document}

\title{Schemes}
\maketitle

\phantomsection
\label{section-phantom}

\tableofcontents

\section{Schemes}
\label{section-shemes}

For a definition of presheaf,
see Categories, Definition
\ref{categories-definition-presheaf}.

\begin{definition}
\label{definition-sheaf}
Let $X$ be a topological space.
\begin{enumerate}
\item A {\it sheaf $\mathcal{F}$ of sets on
$X$} is a presheaf
of sets which satisfies the following
additional property: Given
any open covering $U = \bigcup_{i \in I} U_i$
and any collection
of sections $s_i \in \mathcal{F}(U_i)$, $i
\in I$ such that
$\forall i, j\in I$
$$
s_i|_{U_i \cap U_j} = s_j|_{U_i \cap U_j}
$$
there exists a unique section $s \in
\mathcal{F}(U)$ such that
$s_i = s|_{U_i}$ for all $i \in I$.
\item A {\it morphism of sheaves of sets} is
simply a
morphism of presheaves of sets.
\item The category of sheaves of sets on $X$
is denoted
$\Sh(X)$.
\end{enumerate}
\end{definition}


\medskip\noindent
Let $X$ be a topological space. Let $x \in X$
be a point.
Let $\mathcal{F}$ be a presheaf of sets on
$X$.
The {\it stalk of $\mathcal{F}$ at $x$} is
the set
$$
\mathcal{F}_x
=
\colim_{x\in U} \mathcal{F}(U)
$$
where the colimit is over the set of open
neighbourhoods
$U$ of $x$ in $X$. The set of open
neighbourhoods is
partially ordered by (reverse) inclusion:
We say $U \geq U' \Leftrightarrow U \subset
U'$.
The transition maps in the system are
given by the restriction maps of
$\mathcal{F}$.
See Categories, Section
\ref{categories-section-posets-limits}
for notation and terminology regarding
(co)limits over systems.
Note that the colimit is a directed colimit.
Thus it is easy to describe $\mathcal{F}_x$.
Namely,
$$
\mathcal{F}_x
=
\{
(U, s)
\mid
x\in U, s\in \mathcal{F}(U)
\}/\sim
$$
with equivalence relation given by $(U, s)
\sim (U', s')$ if and only if
there exists an open $U'' \subset U \cap U'$
with $x \in U''$ and
$s|_{U''} = s'|_{U''}$. Given a pair $(U, s)$
we sometimes denote
$s_x$ the element of $\mathcal{F}_x$
corresponding to the equivalence
class of $(U, x)$. We sometimes use the
phrase
``image of $s$ in $\mathcal{F}_x$'' to denote
$s_x$.
For example, given two pairs $(U, s)$ and
$(U', s')$ we sometimes
say ``$s$ is equal to $s'$ in
$\mathcal{F}_x$'' to indicate
that $s_x = s'_x$. Other authors use the
terminology
``germ of $s$ at $x$''.

\section{Closed immersions}
\label{section-immersions}

\noindent
Closed immersions of ringes spaces are in
\href{https://stacks.math.columbia.edu/tag/%
01C1}{01C1}.

\begin{definition}
\label{definition-closed-immersion}
A map of ringed spaces
$i:(Z,\mathcal{O}_Z) \to (X,\mathcal{O}_X)$
is a {\it closed immersion}
if $i$ is injective,
its image is closed (I think these two are
``closed immersion of topological spaces'')
and the induced map $\mathcal{O}_X \to
i_*\mathcal{O}_Z$
is surjective; denote its kernel by
$\mathcal{I}$.
To make sure that if $(X,\mathcal{O}_X)$ is a
scheme
then so is $(Z,\mathcal{O}_Z)$ with add the
extra condition
that the $\mathcal{O}_X$-module $\mathcal{I}$
is
locally generated by sections.
\end{definition}

\medskip\noindent
\section{Affine schemes}
\label{section-affine-schemes}

\noindent
The main objective of this section is
to define the structure sheaf of the
spectrum of a ring, which
was already introduced in 
Commutative Algebra 
\ref{commutative-algebra-section-spectrum}.
The resulting pair
$(\Spec R, \mathcal{O}_{\Spec R}$
will be an affine scheme.

\section{Reduced schemes}
\label{section-reduced}

\noindent
So far a reduced scheme,
for me,
is just a scheme where there are
no ``repetitions'':
$(x^2)$ is not reduced because
it gives the same information as $(x)$.
Algebraically this means that the
coordinate ring is reduced,
which is defined by asking
that there are no nilpotent elements,
i.e. elements such that some power vanishes.
(There's no definition of reduced ring in
Stacks Project
since it's taken as part of the basic
algebraic knowledge).

The definition is on stalks but the first
lemma
shows it's equivalent to define it on any
open set.
This is in virtue of some ``injectivity''
lemma
from the ring at $U$ to the product of all
stalks parametrized in $U$;
essentially that if a section vanishes when
projected to all
the stalks then it's zero.

\begin{definition}
\label{definition-reduced}
Let $X$ be a scheme. We say $X$ is {\it
reduced} if every local ring
$\mathcal{O}_{X, x}$ is reduced.
\end{definition}

\begin{proof}
Assume that $X$ is reduced.
Let $f \in \mathcal{O}_X(U)$ be a section
such that $f^n = 0$.
Then the image of $f$ in $\mathcal{O}_{U, u}$
is zero for
all $u \in U$. Hence $f$ is zero, see
Sheaves, Lemma
\ref{sheaves-lemma-sheaf-subset-stalks}.
Conversely, assume that $\mathcal{O}_X(U)$ is
reduced
for all opens $U$. Pick any nonzero element
$f \in \mathcal{O}_{X, x}$.
Any representative $(U, f \in
\mathcal{O}(U))$  of $f$ is nonzero and
hence not nilpotent. Hence $f$ is not
nilpotent in $\mathcal{O}_{X, x}$.
\end{proof}

\section{Normalization}
\label{section-normalization}

\noindent
It might be interest to check
the construction of
normalization in Stacks Project:
it is not done via gluing.

\begin{definition}
\label{definition-normal-scheme}
\begin{reference}
\cite[Exercise II.3.8]{hart}
\end{reference}
A scheme is {\it normal} 
if all of its local rings are integrally
closed.
\end{definition}

\begin{exercise}[Existence of normalization]
\label{exercise-normalization-schemes}
\begin{reference}
\cite[Exercise II.3.8]{hart}
\end{reference}
Let $X$ be an integral scheme.
For each open affine subset
$U = \Spec A$ of $X$,
let $\tilde{A}$ be the integral closure
of $A$ in its quotient field,
and let $\tilde{U} = \Spec \tilde{A}$.
Show that one can glue the schemes $\tilde{U}$
to obtain a normal integral scheme
$\tilde{X}$, called
the {\it normalization} of $X$.
\end{exercise}

\begin{proof}
The first challenge is to define
the $U_{ij}$. If $U_i = \Spec(\widetilde{A_i})$,
then the most natural is to define
$U_{ij}$ as the spectrum of the 
integral closure of the polynomial ring
modulo the
ideal $I_i + I_j$, which is the ideal
that corresponds
to the intersection of the original
spectrums. This looks like it needs
further reading.
\end{proof}

\medskip\noindent
Here's AI's solution (probably).

Let
$$
\nu : \widetilde C \longrightarrow C
$$
be the normalization of $C$. We recall its
construction. Choose an affine open covering
$U_i=\Spec(A_i)$ of $C$. Since $C$ is
irreducible, all $A_i$ have the same
function field $K(C)$. Let $\overline A_i$
be the integral closure of $A_i$ in $K(C)$.

If $U_{ij}=U_i\cap U_j$, then, after
refining the cover if necessary, we may
assume
$$
U_{ij}=D(f)\subset U_i
$$
for some $f\in A_i$. Its coordinate ring is
$(A_i)_f$. The integral closure of
$(A_i)_f$ in $K(C)$ is
$$
(\overline A_i)_f,
$$
because integral closure commutes with
localization. Thus the normalization over
$U_{ij}$ is
$$
\Spec((\overline A_i)_f).
$$
Doing the same computation from $U_j$ gives
the same subring of $K(C)$, namely the
integral closure of $\mathcal O_C(U_{ij})$
inside $K(C)$. Hence the affine schemes
$\Spec(\overline A_i)$ have canonical
identifications over the intersections
$U_i\cap U_j$. These identifications satisfy
the cocycle condition, since they are all
induced by equality inside $K(C)$.

Therefore the affine schemes
$\Spec(\overline A_i)$ glue to a curve
$\widetilde C$, and the inclusions
$$
A_i\subset \overline A_i
$$
glue to a finite birational morphism
$$
\nu:\widetilde C\to C.
$$
The curve $\widetilde C$ is normal by
construction.


\section{Proj}
\label{section-proj}

\noindent
(Stacks Project
\href{https://stacks.math.columbia.edu/tag/%
01M3}{01M3},
Proj of a graded ring.)
This section may be taken as the preamble to
the question: what is projective space
anyway?
Yes, projective space is Proj of a graded
ring $S$.
And yes, projective space is a scheme.
So in particular it's locally affine.

(Stacks Project
\href{https://stacks.math.columbia.edu/tag/%
01MX}{01MX},
Functoriality of Proj.)
This is reminiscent of the olden days.
The point is that Proj is covariant,
i.e. for rings  $S' \subset S$ we have
$\text{Proj}S \subset \text{Proj}S$
(under the right conditions,
which is the point of this section),
meaning that, indeed,
{\bf Projs are projective schemes}.






\medskip\noindent
Noe let's discuss invertible sheaves on Proj
(twists!).

\noindent
As for the construction of the twisted
sheaves:
for a projective space $\text{Proj} S$,
we just compute the twisted $S$ module
$S(m)$ by the formula $S(m)_d=S_{m+d}$
and turn that into a sheaf
$\widetilde{S(m)}$.

Here's the construction of the twisted
sheaves $\mathcal{O}_X(n)$.
The point is that there is a good way
to pass from an $S$-module $M$ to an
$\mathcal{O}_X$-module, called
$\widetilde{M}$.
So basically you just define the twists
by $M(d)_n=M_{n+d}$ and apply that
construction.

Pick an element of degree $d$.
Then think of the module generated by this
element.
The least degree piece of such a module
is the degree $d$ piece of the original
module,
that is we have $M(-d)$.










\section{Restricting sheaves to subschemes}
\label{section-restricting-sheaves}

\noindent
Upshot about restricting sheaves to
subschemes:
the sheaves on $Z$ correspond essentially
(in the categorical sense)
to sheaves on $X$ with support on $Z$.

More explicitly, from
\href{https://stacks.math.columbia.edu/tag/%
01QY}{01QY}:
the pushforward of the inclusion (a closed
immersion of schemes)
$i:Z \to X$, with kernel on induced map
$\mathcal{O}_Z \to \mathcal{O}_X$
denoted by $I$, is an exact, fully faithful
functor
$$
i_*:\QCoh(\mathcal{O}_Z) \to
\QCoh(\mathcal{O}_X)
$$
with essential image the quasi-coherent
$\mathcal{O}_X$ modules
$\mathcal{G}$ such that
$\mathcal{I}\mathcal{G}=0$.

It makes sense to think of this operation as
$\mathcal{F}|_Z=i^*\mathcal{F}=\mathcal{F}
\otimes\mathcal{O}_Z$
for $\mathcal{F}\in\QCoh(X)$.
Indeed: tensoring by the ideal sheaf would
kill the information on $Z$,
and tensoring with the structure sheaf does
the opposite
--- it's the sheaves on $X$ that are no fun
outside $Z$,
so basically sheaves on $Z$.

\medskip\noindent
Here's a more detailed discussion:

In understanding the cohomology of sheaves
when
we restrict them to submanifolds (see Complex
Geometry
Exercise
\ref{complex-geometry-exercise-%
deformations-curve-in-k3},
I found Stacks Project
\href{https://stacks.math.columbia.edu/tag/%
01AW}{01AW}
section on closed immersions and abelian
sheaves.
First we have a few propositions stating that
sheaves on the subscheme $Z \subset X$ can be
pushed
to sheaves on $X$.

And then there's the converse:
what I would like to call the restriction
functor $\mathcal{H}$
but rather is called the {\it sub-module of
sections with support in $Z$}.
In Remark
\href{https://stacks.math.columbia.edu/tag/%
01AY}{01AY}
(and then in Remark
\href{https://stacks.math.columbia.edu/tag/%
0G6N}{0G6N}
for ringed spaces; these seem to be analogue
constructions…)
we see that for an abelian sheaf
$\mathcal{F}$ on $X$
we can define an abelian sheaf
$\mathcal{H}_Z(\mathcal{F})$
which although is a sheaf on $X$
we can just think it's a sheaf on $Z$ -
that's beacause
$$
\mathcal{H}_Z(\mathcal{F})(U)
=\{s \in \mathcal{F}(U))| \text{the support
of $s$ is contained in }Z\cap U\}.
$$
We can just think that it's the functions
that vanish outside the subvariety.
Which doesn't really make sense:
a function vanishing outside the submanifold
will most likely
vanish also in the submanifold! Just by
continuity,
the submanifold is in the closure of the open
set!
I guess what I mean to say is that I don't
mind if
the functions do not vanish outside $Z$…
I just want to consider them as functions on
$Z$…

\medskip\noindent
Here's another attempt in understanding how
to restrict sheaves to submanifolds.
First and foremost, lacking a reference this,
I understand the symbol $\mathcal{F}|_Z$ for
a sheaf
on $X$ and a subscheme  $Z \subset X$ to mean
the
pullback of $\mathcal{F}$ by the inclusion.
Then look at
\href{https://stacks.math.columbia.edu/tag/%
01QY}{01QY}
(more basically, as above,
\href{https://stacks.math.columbia.edu/tag/%
01AW}{01AW}…
in fact there appear to be several version of
this statement)
to recall that sheaves of the subscheme $Z$
should correspond
to sheaves on $X$ with support on $Z$.
{\bf So essentially $\mathcal{F}|_Z$ should
be a version of $\mathcal{F}$
but with support on $Z$.}
The key point here is that this is non other
than
$\mathcal{F} \otimes \mathcal{O}_Z$.
Why? Because when multiply with the ideal
sheaf of $Z$ we obtain
$(\mathcal{F} \otimes
\mathcal{O}_Z)\mathcal{I}=0$
since sections of $\mathcal{I}$
vanish in
$\mathcal{O}_Z=\mathcal{O}_X/\mathcal{I}$.
That is, the support of $\mathcal{F} \otimes
\mathcal{O}_Z$ indeed
contained in $Z$: if the sections become zero
when
multiplied by sections of $\mathcal{I}$
(sections of $\mathcal{I}$ are only nonzero
outside $Z$)
it means that all the action was inside $Z$.
that outside $Z$

I think the upshot here is that structure
sheaves
act as some sort of ``contour'' for the
geometric object:
they vanish {\it outside} the geometric
object,
they tell us what is the shape of the object
by describing its surroundings, rather than
its interior.
This is exactly how algebraic/analytic sets
are defined,
and this why, when we want to restrict a
sheaf
to a subscheme, we tensor with the structure
sheaf:
tensor by the ideal sheaf would kill the
information on $Z$,
and tensoring with the structure sheaf does
the opposite.













\section{Flat modules}
\label{section-flat-modules}

\noindent
The essential technical property
for for defining flatness
is the preservation of exact sequences.
Right-exactness is true in general;
it follows from currying in category
of commutative rings.
But the functor $-\otimes_R N$ does not
preserve injectivity of maps---
that's the point of flatness.

\begin{lemma}[Internal Hom for R-modules]
\label{lemma-hom-from-tensor-product}
For any three $R$-modules $M, N, P$,
$$
\Hom_R(M \otimes_R N, P) \cong \Hom_R(M,
\Hom_R(N, P))
$$
\end{lemma}

\begin{proof}
An $R$-linear map $\hat{f}\in \Hom_R(M
\otimes_R N, P)$ corresponds to an
$R$-bilinear map $f : M \times N \to P$. For
each $x\in M$ the mapping $y\mapsto f(x, y)$
is $R$-linear by the universal
property. Thus $f$ corresponds to a
map $\phi_f : M \to \Hom_R(N, P)$. This map
is $R$-linear since
$$
\phi_f(ax + y)(z) =
f(ax + y, z) = af(x, z)+f(y, z) =
(a\phi_f(x)+\phi_f(y))(z),
$$
for all $a \in R$, $x \in M$, $y \in M$ and
$z \in N$. Conversely, any
$f \in \Hom_R(M, \Hom_R(N, P))$ defines an
$R$-bilinear
map $M \times N \to P$, namely $(x, y)\mapsto
f(x)(y)$.
So this is a natural one-to-one
correspondence between the
two modules
$\Hom_R(M \otimes_R N, P)$ and $\Hom_R(M,
\Hom_R(N, P))$.
\end{proof}

\begin{lemma}[Tensor product is right exact]
\label{lemma-tensor-product-exact}
Let
$$
\begin{aligned}
M_1\xrightarrow{f} M_2\xrightarrow{g} M_3 \to
0
\end{aligned}
$$
be an exact sequence of $R$-modules and
homomorphisms, and let $N$ be any
$R$-module. Then the sequence
\begin{equation}
\label{equation-2ndex}
M_1\otimes N\xrightarrow{f \otimes 1}
M_2\otimes N \xrightarrow{g \otimes 1}
M_3\otimes N \to 0
\end{equation}
is exact. In other words, the functor $-
\otimes_R N$ is
{\it right exact}, in the sense that
tensoring
each term in the original right exact
sequence preserves the exactness.
\end{lemma}

\begin{proof}
For every $R$-module $P$
we apply the functor $\Hom(-, \Hom(N, P))$ to
the first exact
sequence. We obtain
$$
0 \to
\Hom(M_3, \Hom(N, P)) \to
\Hom(M_2, \Hom(N, P)) \to
\Hom(M_1, \Hom(N, P))
$$
which is exact by Lemma \ref{lemma-hom-exact}
(1).
By Lemma \ref{lemma-hom-from-tensor-product}
this becomes the sequence
$$
0 \to \Hom(M_3 \otimes N, P) \to
\Hom(M_2 \otimes N, P) \to \Hom(M_1 \otimes
N, P)
$$
which is therefore also exact. Then using
Lemma \ref{lemma-hom-exact} (1) again, we
arrive at the desired exact sequence.
\end{proof}

\begin{remark}
\label{remark-tensor-product-not-exact}
However, tensor product does NOT preserve
exact sequences in general.
In other words, if $M_1 \to M_2 \to M_3$ is
exact, then it is not necessarily true that
$M_1 \otimes N \to M_2 \otimes N \to M_3
\otimes N$
is exact for arbitrary $R$-module $N$.
\end{remark}

\begin{example}
\label{example-tensor-product-not-exact}
Consider the injective map $2 : \mathbf{Z}\to
\mathbf{Z}$
viewed as a map of $\mathbf{Z}$-modules.
Let $N = \mathbf{Z}/2$. Then the induced map
$\mathbf{Z} \otimes \mathbf{Z}/2 \to
\mathbf{Z} \otimes \mathbf{Z}/2$
is NOT injective. This is because for
$x \otimes y\in \mathbf{Z} \otimes
\mathbf{Z}/2$,
$$
(2 \otimes 1)(x \otimes y) = 2x \otimes y = x
\otimes 2y = x \otimes 0 = 0
$$
Therefore the induced map is the zero map
while $\mathbf{Z} \otimes N\neq 0$.
\end{example}

\begin{definition}
\label{definition-flat-module}
For $R$-modules $N$, if the
functor $-\otimes_R N$ is exact, i.e.
tensoring
with $N$ preserves all exact
sequences, then $N$ is said to be {\it flat}
$R$-module.
We will discuss this later in Section
\ref{section-flat}.
\end{definition}

\medskip\noindent
Epiphany:
in the category of commutative rings
pushout is tensor product.
So think of $\Spec $ as a functor
from $\mathit{CRing}^{\text{op}}$ to $\Sch$,
then pushout goes to pullback,
and what's an example of a pullback?
Fibre!
So, the coordinate ring of a fiber
is essentially given by the residue field
at the point that parametrizes it!
(tensored with the coordinate ring
of the deformation space).

There is a lot of information on Stacks
Project about flatness.
It looks like the heart of the concept is
captured in the commutative-algebraic notion
of preserving
exact sequences:

\begin{definition}
\label{definition-flat}
Let $R$ be a ring.
\begin{enumerate}
\item An $R$-module $M$ is called {\it flat}
if whenever
$N_1 \to N_2 \to N_3$ is an exact sequence of
$R$-modules
the sequence $M \otimes_R N_1 \to M \otimes_R
N_2 \to M \otimes_R N_3$
is exact as well.
\item An $R$-module $M$ is called {\it
faithfully flat} if the
complex of $R$-modules
$N_1 \to N_2 \to N_3$ is exact if and only if
the sequence $M \otimes_R N_1 \to M \otimes_R
N_2 \to M \otimes_R N_3$
is exact.
\item A ring map $R \to S$ is called {\it
flat} if
$S$ is flat as an $R$-module.
\item A ring map $R \to S$ is called {\it
faithfully flat} if
$S$ is faithfully flat as an $R$-module.
\end{enumerate}
\end{definition}

\medskip\noindent
Recall that a module $M$ over a ring $R$ is
{\it flat} if the functor
$-\otimes_R M : \text{Mod}_R \to
\text{Mod}_R$ is exact. A ring map
$R \to A$ is said to be {\it flat} if $A$ is
flat as an $R$-module.


\bibliography{my}
\bibliographystyle{amsalpha}

\end{document}

